\documentclass{article}
\usepackage[UTF8]{ctex}
\usepackage{amsmath,amssymb,amsfonts}
\usepackage{enumitem}
\usepackage{geometry}
\geometry{a4paper, margin=1in}
% === 错题本排版优化（不参与题目标识） ===
\setlength{\parindent}{0pt}
\setlength{\parskip}{0.35em plus 0.10em minus 0.08em}
\setlength{\emergencystretch}{3em}
\setlist{itemsep=0.25em, topsep=0.40em, parsep=0pt}
\allowdisplaybreaks[2]
% === 错题本排版优化结束 ===
\begin{document}
\section*{东南大学2025年数学分析真题}
\section*{87．东南大学2025年数学分析试卷（回忆版）}
\begin{itemize}

  \item[1．] 求空间曲线
\[
\left\{\begin{array}{l}
x+y+z=0 \cdots \\
x^{2}+y^{2}=z^{2} \cdots
\end{array}\right.
\]
的切线，其满足平行于平面 \(x+2 y=0 \cdots\) ．

  \item[2．] \(z=z(x, y)\) 是由方程 \(\cdots-\int_{y}^{-\sqrt{x^{2}+y^{2}}} e^{-t^{2}} \mathrm{~d} t-\cdots=0\) 确定，求 \(\frac{\partial z}{\partial x}, \frac{\partial^{2} z}{\partial x \partial y}\) ．
\end{itemize}

3．求幂级数 \(\sum_{n=1}^{\infty}(-1)^{n-1} n^{2}(x-1)^{n}\) 的收敛域与和函数。

4．求函数 \(f(x)=\left\{\begin{array}{ll}\frac{1}{x \ln 2}-\frac{1}{2^{x}-1}, & x \neq 0 ; \\ \frac{1}{2}, & x=0 .\end{array}\right.\) 在 \(x=0\) 的导数．

\begin{itemize}
\end{itemize}

\subsection{题目5}
原题边界“5．”：设 \(f(x)\) 在 \([a, b]\) 上连续，在 \((a, b)\) 上可导，且 \(f(x) \neq 0\) ，证明：存在 \(\xi \in(a, b)\) ，满足
\[
\frac{f^{\prime}(\xi)}{f(\xi)}=\frac{1}{a-\xi}+\frac{1}{b-\xi} \ldots
\]

\subsection{题目6}
原题边界“6．”：求 \(\int_{L} \frac{x \mathrm{~d} y-y \mathrm{~d} x}{x^{2}+4 y^{2}}, L\) 是曲线 \((x-1)^{2}+y^{2}=r^{2}(r \neq 1)\) ，取正向．

7．求 \(f(x, y, z)=x^{3}+y^{3}+z^{3}-2 x y z\) 的最值，其中 \(x^{2}+y^{2}+z^{2}=4\) 。

8．求 \(\iiint_{\Omega}|x y| z \mathrm{~d} x \mathrm{~d} y \mathrm{~d} z\) ，其中 \(\Omega\) 是由 \(x+z=1, x-z=1, x+y=1, x-y=1, z=0\) 围成的区域。

9．设 \(F(x)=\int_{0}^{+\infty} e^{-y^{2}} \cos 2 x y \mathrm{~d} y\) ．

\begin{itemize}
  \item[（1）] 证明： \(2 F^{\prime}(x)+x F(x)=0\) ．
  \item[（2）] 求 \(F(x)\) 。
\end{itemize}

10．设 \(f(x)\) 在 \([a,+\infty)\) 上一致连续，\(g(x)\) 在 \([a,+\infty)\) 上连续，且 \(\lim _{x \rightarrow+\infty}[f(x)-g(x)]=0\) ，证明：\(g(x)\)在 \([a,+\infty)\) 上一致连续。

11．求 \(\iint_{\sum} z \mathrm{~d} S\) ，其中 \(\sum\) 是由 \(x^{2}+z^{2}=2 x\) 与 \(z=\sqrt{x^{2}+y^{2}}\) 为围成的有限曲面．

\begin{itemize}
\end{itemize}

\subsection{题目12}
原题边界“12．”：记 \(u_{n}(x)=\frac{\ln \left(1+n^{\alpha} x\right)}{n^{\beta}}, x>0, \alpha, \beta>0\) ．

\begin{itemize}
\item[（1）] 讨论 \(\sum_{n=1}^{\infty} u_{n}(x)\) 的敛散性．
  \item[（2）] 在 \(\sum_{n=1}^{\infty} u_{n}(x)\) 收敛的基础上，讨论其一致收敛性。
  \item[（3）] 讨论 \(\sum_{n=1}^{\infty} u_{n}(x)\) 的可导性．
\end{itemize}

13．设 \(f: \mathbb{R} \rightarrow \mathbb{R}\) 为连续函数，\(f\) 无不动点，记

\[
f_{n}(x)=\underbrace{f \circ f \circ f \circ \cdots \circ f}_{n \text { 个 }}(x)
\]

为 \(f(x)\) 的 \(n\) 次复合，数列 \(\left\{f_{n}(x)\right\}\) 是否有界？
\end{document}
